\documentclass[../../main.tex]{subfiles}
\begin{document}

\begin{flushright}
  \footnotesize\textit{【自動文字起こし・要確認】}
\end{flushright}

% Source: handwritten.pdf 第3問．
{\bf [解]}
与式から
\begin{align}
  a+d=-(b+c),\label{eq:ky2007-3-a}\\
  ad-bc=-p,\label{eq:ky2007-3-b}\\
  d\leqq c\leqq b\leqq a.\label{eq:ky2007-3-c}
\end{align}
$a+d=\alpha$，$ad=\beta$とおく．$\alpha,\beta\in\mathbb{Z}$で，\cref{eq:ky2007-3-a,eq:ky2007-3-b}から，
$a,d$は$X^2-\alpha X+\beta=0$の2解，$b,c$は$X^2+\alpha X+\beta+p=0$の2解である．
だから
\begin{align}
  \begin{cases}
    a^2-\alpha a+\beta=0,\\
    d^2-\alpha d+\beta=0,\\
    c^2+\alpha c+\beta+p=0,\\
    b^2+\alpha b+\beta+p=0.
  \end{cases}\label{eq:ky2007-3-roots}
\end{align}
$a$と$b$の式を辺々引いて，
\begin{align}
  a^2-b^2-\alpha(a+b)=p,\qquad(a+b)(a-b-\alpha)=p,
\end{align}
\begin{align}-(a+b)(b+d)=p.
\end{align}
ここで$a,b,c,d\in\mathbb{Z}$，$p$は素数であり，\cref{eq:ky2007-3-c}から$b+d\leqq a+b$であることより，
\begin{align}(b+d,a+b)=(-1,p),(-p,1).\label{eq:ky2007-3-d}
\end{align}
同じことを$a$と$c$の式を用いて行えば，
\begin{align}-(a+c)(c+d)=p
\end{align}
から
\begin{align}(c+d,a+c)=(-1,p),(-p,1).\label{eq:ky2007-3-e}
\end{align}
\cref{eq:ky2007-3-c,eq:ky2007-3-d,eq:ky2007-3-e}及び
$c+d\leqq b+d\leqq a+c\leqq a+b$だから，
\begin{align}
  (c+d,b+d,a+c,a+b)=(-1,-1,p,p),\quad(-p,-1,1,p),\quad(-p,-p,1,1)
\end{align}
である．

\subparagraph{(i) $(-1,-1,p,p)$のとき}
$c=-1-d=p-a$，$b=-1-d=p-a$となるから，$a-d=p+1$となる．
\cref{eq:ky2007-3-a,eq:ky2007-3-b,eq:ky2007-3-c}に代入して，
\begin{align}
  \begin{cases}
    (d+p+1)+d=2(1+d),\\
    d(d+p+1)-(1+d)^2=-p,\\
    d\leqq-1-d=-1-d\leqq d+p+1.
  \end{cases}
\end{align}
第1式から$p=1$となって矛盾．

\subparagraph{(ii) $(-p,-1,1,p)$のとき}
$c=-p-d=1-a$，$b=-1-d=p-a$から，$a=d+p+1$とする．
\cref{eq:ky2007-3-a,eq:ky2007-3-b,eq:ky2007-3-c}に代入して，
\begin{align}
  \begin{cases}
    2d+p+1=2d+p+1,\\
    d(d+p+1)-(p+d)(1+d)=-p,\\
    d\leqq-p-d\leqq-1-d\leqq d+p+1.
  \end{cases}
\end{align}
したがって$d$を$-p-2\leqq2d\leqq-p$であればよく，つまり
\begin{align}d=-\dfrac{p+1}{2}
\end{align}
とすればよい．$p$は3以上の素数であるから，このとき全ての条件がみたされる．
\begin{align}
  a=\dfrac{p+1}{2},\qquad b=\dfrac{p-1}{2},\qquad c=-\dfrac{p-1}{2}.
\end{align}

\subparagraph{(iii) $(-p,-p,1,1)$のとき}
$c=-p-d=1-a$，$b=-p-d=1-a$から，$a=p+d+1$となる．
\cref{eq:ky2007-3-a}に代入して
\begin{align}2d+p+1=2p+2d
\end{align}
から$p=1$となって矛盾．
\casesend
以上から
\begin{align}
  (a,b,c,d)=\left(\dfrac{p+1}{2},\dfrac{p-1}{2},-\dfrac{p-1}{2},-\dfrac{p+1}{2}\right).
\end{align}

\medskip\noindent{\bf [補足]}
\cref{eq:ky2007-3-roots}の$a$と$b$の式において，与式から$a+b=-(c+d)$だから，
\cref{eq:ky2007-3-d}のみに注目する方がよい．

\medskip\noindent{\bf [別解についての原稿メモ]}
普通に文字消去でもとける．というか，その方が早かった．

\end{document}
